% Business Model Canvas. Una página horizontal dentro de modelo_negocio.tex, que
% lo inserta con % Business Model Canvas. Una página horizontal dentro de modelo_negocio.tex, que
% lo inserta con % Business Model Canvas. Una página horizontal dentro de modelo_negocio.tex, que
% lo inserta con % Business Model Canvas. Una página horizontal dentro de modelo_negocio.tex, que
% lo inserta con \input{canvas}. Este archivo no se compila solo.
% Llenar los nueve bloques. Cada afirmación termina con el identificador de la
% hipótesis que la respalda en hipotesis.csv, por ejemplo (H1).
% Marcar las hipótesis validadas con \val{H1} y las invalidadas con \inval{H2}.
\definecolor{navy}{HTML}{1F3467}
\providecommand{\val}[1]{\textcolor{green!50!black}{\textbf{(#1 \checkmark)}}}
\providecommand{\inval}[1]{\textcolor{red!70!black}{\textbf{(#1 $\times$)}}}
\providecommand{\pendiente}[1]{\textcolor{gray}{\emph{#1}}}

% ---------- Contenido de los bloques ----------
\newcommand{\Socios}{\begin{itemize}
  \item \pendiente{Proveedores y aliados sin los que el negocio no opera.} (H?)
\end{itemize}}
\newcommand{\Actividades}{\begin{itemize}
  \item \pendiente{Lo que el equipo hace todos los días para entregar valor.} (H?)
\end{itemize}}
\newcommand{\Recursos}{\begin{itemize}
  \item \pendiente{Datos, modelos, personas e infraestructura.} (H?)
\end{itemize}}
\newcommand{\Propuesta}{\begin{itemize}
  \item \pendiente{Qué dolor elimina o qué ganancia crea, para quién.} (H?)
\end{itemize}}
\newcommand{\Relacion}{\begin{itemize}
  \item \pendiente{Autoservicio, asistencia personal, comunidad.} (H?)
\end{itemize}}
\newcommand{\Canales}{\begin{itemize}
  \item \pendiente{Cómo se entera, compra y recibe el producto el cliente.} (H?)
\end{itemize}}
\newcommand{\Segmentos}{\begin{itemize}
  \item \pendiente{Perfil del cliente que paga y del usuario, si son distintos.} (H?)
\end{itemize}}
\newcommand{\Costos}{\begin{itemize}
  \item \pendiente{Costos fijos y variables principales (ver modelo financiero).} (H?)
\end{itemize}}
\newcommand{\Ingresos}{\begin{itemize}
  \item \pendiente{Qué se cobra, cuánto y con qué frecuencia.} (H?)
\end{itemize}}

% ---------- Dibujo. No editar debajo de esta línea ----------
\begingroup
\fontsize{9}{11}\selectfont
\setlist{nosep,leftmargin=1em}
\noindent{\Large\bfseries\color{navy} Business Model Canvas}\hfill
{\large \pendiente{Nombre de la empresa}}\hfill \pendiente{Versión y fecha}

\vspace{4pt}
\noindent
\resizebox{\linewidth}{!}{%
\begin{tikzpicture}[x=1cm,y=1cm]
  \def\W{25.8}\def\H{16.8}\def\h{5.0}
  \tikzset{caja/.style={draw=navy,line width=0.8pt},
           tit/.style={anchor=north west,font=\bfseries\color{navy}},
           cue/.style={anchor=north west,text width=#1cm,align=left}}
  % Retícula: 5 columnas arriba, 2 abajo
  \def\c{5.16}
  \draw[caja] (0,0) rectangle (\W,\H);
  \draw[caja] (0,\h) -- (\W,\h);
  \foreach \i in {1,2,3,4} \draw[caja] (\i*\c,\h) -- (\i*\c,\H);
  \draw[caja] (\W/2,0) -- (\W/2,\h);
  \pgfmathsetmacro\m{(\h+\H)/2}
  \draw[caja] (\c,\m) -- (2*\c,\m);
  \draw[caja] (3*\c,\m) -- (4*\c,\m);
  % Títulos y contenido
  \node[tit] at (0.1,\H-0.1) {Socios clave};
  \node[cue=4.9] at (0.1,\H-0.7) {\Socios};
  \node[tit] at (\c+0.1,\H-0.1) {Actividades clave};
  \node[cue=4.9] at (\c+0.1,\H-0.7) {\Actividades};
  \node[tit] at (\c+0.1,\m-0.1) {Recursos clave};
  \node[cue=4.9] at (\c+0.1,\m-0.7) {\Recursos};
  \node[tit] at (2*\c+0.1,\H-0.1) {Propuesta de valor};
  \node[cue=4.9] at (2*\c+0.1,\H-0.7) {\Propuesta};
  \node[tit] at (3*\c+0.1,\H-0.1) {Relación con clientes};
  \node[cue=4.9] at (3*\c+0.1,\H-0.7) {\Relacion};
  \node[tit] at (3*\c+0.1,\m-0.1) {Canales};
  \node[cue=4.9] at (3*\c+0.1,\m-0.7) {\Canales};
  \node[tit] at (4*\c+0.1,\H-0.1) {Segmentos de clientes};
  \node[cue=4.9] at (4*\c+0.1,\H-0.7) {\Segmentos};
  \node[tit] at (0.1,\h-0.1) {Estructura de costos};
  \node[cue=12.6] at (0.1,\h-0.7) {\Costos};
  \node[tit] at (\W/2+0.1,\h-0.1) {Fuentes de ingresos};
  \node[cue=12.6] at (\W/2+0.1,\h-0.7) {\Ingresos};
\end{tikzpicture}}

\vspace{2pt}
\noindent\scriptsize Basado en Osterwalder y Pigneur, \emph{Business Model Generation} (2010), CC BY-SA 3.0.
\val{H1} validada,\ \inval{H2} invalidada, (H3) pendiente. Evidencia en \texttt{hipotesis.csv} y \texttt{entrevistas.csv}.
\par\endgroup
. Este archivo no se compila solo.
% Llenar los nueve bloques. Cada afirmación termina con el identificador de la
% hipótesis que la respalda en hipotesis.csv, por ejemplo (H1).
% Marcar las hipótesis validadas con \val{H1} y las invalidadas con \inval{H2}.
\definecolor{navy}{HTML}{1F3467}
\providecommand{\val}[1]{\textcolor{green!50!black}{\textbf{(#1 \checkmark)}}}
\providecommand{\inval}[1]{\textcolor{red!70!black}{\textbf{(#1 $\times$)}}}
\providecommand{\pendiente}[1]{\textcolor{gray}{\emph{#1}}}

% ---------- Contenido de los bloques ----------
\newcommand{\Socios}{\begin{itemize}
  \item \pendiente{Proveedores y aliados sin los que el negocio no opera.} (H?)
\end{itemize}}
\newcommand{\Actividades}{\begin{itemize}
  \item \pendiente{Lo que el equipo hace todos los días para entregar valor.} (H?)
\end{itemize}}
\newcommand{\Recursos}{\begin{itemize}
  \item \pendiente{Datos, modelos, personas e infraestructura.} (H?)
\end{itemize}}
\newcommand{\Propuesta}{\begin{itemize}
  \item \pendiente{Qué dolor elimina o qué ganancia crea, para quién.} (H?)
\end{itemize}}
\newcommand{\Relacion}{\begin{itemize}
  \item \pendiente{Autoservicio, asistencia personal, comunidad.} (H?)
\end{itemize}}
\newcommand{\Canales}{\begin{itemize}
  \item \pendiente{Cómo se entera, compra y recibe el producto el cliente.} (H?)
\end{itemize}}
\newcommand{\Segmentos}{\begin{itemize}
  \item \pendiente{Perfil del cliente que paga y del usuario, si son distintos.} (H?)
\end{itemize}}
\newcommand{\Costos}{\begin{itemize}
  \item \pendiente{Costos fijos y variables principales (ver modelo financiero).} (H?)
\end{itemize}}
\newcommand{\Ingresos}{\begin{itemize}
  \item \pendiente{Qué se cobra, cuánto y con qué frecuencia.} (H?)
\end{itemize}}

% ---------- Dibujo. No editar debajo de esta línea ----------
\begingroup
\fontsize{9}{11}\selectfont
\setlist{nosep,leftmargin=1em}
\noindent{\Large\bfseries\color{navy} Business Model Canvas}\hfill
{\large \pendiente{Nombre de la empresa}}\hfill \pendiente{Versión y fecha}

\vspace{4pt}
\noindent
\resizebox{\linewidth}{!}{%
\begin{tikzpicture}[x=1cm,y=1cm]
  \def\W{25.8}\def\H{16.8}\def\h{5.0}
  \tikzset{caja/.style={draw=navy,line width=0.8pt},
           tit/.style={anchor=north west,font=\bfseries\color{navy}},
           cue/.style={anchor=north west,text width=#1cm,align=left}}
  % Retícula: 5 columnas arriba, 2 abajo
  \def\c{5.16}
  \draw[caja] (0,0) rectangle (\W,\H);
  \draw[caja] (0,\h) -- (\W,\h);
  \foreach \i in {1,2,3,4} \draw[caja] (\i*\c,\h) -- (\i*\c,\H);
  \draw[caja] (\W/2,0) -- (\W/2,\h);
  \pgfmathsetmacro\m{(\h+\H)/2}
  \draw[caja] (\c,\m) -- (2*\c,\m);
  \draw[caja] (3*\c,\m) -- (4*\c,\m);
  % Títulos y contenido
  \node[tit] at (0.1,\H-0.1) {Socios clave};
  \node[cue=4.9] at (0.1,\H-0.7) {\Socios};
  \node[tit] at (\c+0.1,\H-0.1) {Actividades clave};
  \node[cue=4.9] at (\c+0.1,\H-0.7) {\Actividades};
  \node[tit] at (\c+0.1,\m-0.1) {Recursos clave};
  \node[cue=4.9] at (\c+0.1,\m-0.7) {\Recursos};
  \node[tit] at (2*\c+0.1,\H-0.1) {Propuesta de valor};
  \node[cue=4.9] at (2*\c+0.1,\H-0.7) {\Propuesta};
  \node[tit] at (3*\c+0.1,\H-0.1) {Relación con clientes};
  \node[cue=4.9] at (3*\c+0.1,\H-0.7) {\Relacion};
  \node[tit] at (3*\c+0.1,\m-0.1) {Canales};
  \node[cue=4.9] at (3*\c+0.1,\m-0.7) {\Canales};
  \node[tit] at (4*\c+0.1,\H-0.1) {Segmentos de clientes};
  \node[cue=4.9] at (4*\c+0.1,\H-0.7) {\Segmentos};
  \node[tit] at (0.1,\h-0.1) {Estructura de costos};
  \node[cue=12.6] at (0.1,\h-0.7) {\Costos};
  \node[tit] at (\W/2+0.1,\h-0.1) {Fuentes de ingresos};
  \node[cue=12.6] at (\W/2+0.1,\h-0.7) {\Ingresos};
\end{tikzpicture}}

\vspace{2pt}
\noindent\scriptsize Basado en Osterwalder y Pigneur, \emph{Business Model Generation} (2010), CC BY-SA 3.0.
\val{H1} validada,\ \inval{H2} invalidada, (H3) pendiente. Evidencia en \texttt{hipotesis.csv} y \texttt{entrevistas.csv}.
\par\endgroup
. Este archivo no se compila solo.
% Llenar los nueve bloques. Cada afirmación termina con el identificador de la
% hipótesis que la respalda en hipotesis.csv, por ejemplo (H1).
% Marcar las hipótesis validadas con \val{H1} y las invalidadas con \inval{H2}.
\definecolor{navy}{HTML}{1F3467}
\providecommand{\val}[1]{\textcolor{green!50!black}{\textbf{(#1 \checkmark)}}}
\providecommand{\inval}[1]{\textcolor{red!70!black}{\textbf{(#1 $\times$)}}}
\providecommand{\pendiente}[1]{\textcolor{gray}{\emph{#1}}}

% ---------- Contenido de los bloques ----------
\newcommand{\Socios}{\begin{itemize}
  \item \pendiente{Proveedores y aliados sin los que el negocio no opera.} (H?)
\end{itemize}}
\newcommand{\Actividades}{\begin{itemize}
  \item \pendiente{Lo que el equipo hace todos los días para entregar valor.} (H?)
\end{itemize}}
\newcommand{\Recursos}{\begin{itemize}
  \item \pendiente{Datos, modelos, personas e infraestructura.} (H?)
\end{itemize}}
\newcommand{\Propuesta}{\begin{itemize}
  \item \pendiente{Qué dolor elimina o qué ganancia crea, para quién.} (H?)
\end{itemize}}
\newcommand{\Relacion}{\begin{itemize}
  \item \pendiente{Autoservicio, asistencia personal, comunidad.} (H?)
\end{itemize}}
\newcommand{\Canales}{\begin{itemize}
  \item \pendiente{Cómo se entera, compra y recibe el producto el cliente.} (H?)
\end{itemize}}
\newcommand{\Segmentos}{\begin{itemize}
  \item \pendiente{Perfil del cliente que paga y del usuario, si son distintos.} (H?)
\end{itemize}}
\newcommand{\Costos}{\begin{itemize}
  \item \pendiente{Costos fijos y variables principales (ver modelo financiero).} (H?)
\end{itemize}}
\newcommand{\Ingresos}{\begin{itemize}
  \item \pendiente{Qué se cobra, cuánto y con qué frecuencia.} (H?)
\end{itemize}}

% ---------- Dibujo. No editar debajo de esta línea ----------
\begingroup
\fontsize{9}{11}\selectfont
\setlist{nosep,leftmargin=1em}
\noindent{\Large\bfseries\color{navy} Business Model Canvas}\hfill
{\large \pendiente{Nombre de la empresa}}\hfill \pendiente{Versión y fecha}

\vspace{4pt}
\noindent
\resizebox{\linewidth}{!}{%
\begin{tikzpicture}[x=1cm,y=1cm]
  \def\W{25.8}\def\H{16.8}\def\h{5.0}
  \tikzset{caja/.style={draw=navy,line width=0.8pt},
           tit/.style={anchor=north west,font=\bfseries\color{navy}},
           cue/.style={anchor=north west,text width=#1cm,align=left}}
  % Retícula: 5 columnas arriba, 2 abajo
  \def\c{5.16}
  \draw[caja] (0,0) rectangle (\W,\H);
  \draw[caja] (0,\h) -- (\W,\h);
  \foreach \i in {1,2,3,4} \draw[caja] (\i*\c,\h) -- (\i*\c,\H);
  \draw[caja] (\W/2,0) -- (\W/2,\h);
  \pgfmathsetmacro\m{(\h+\H)/2}
  \draw[caja] (\c,\m) -- (2*\c,\m);
  \draw[caja] (3*\c,\m) -- (4*\c,\m);
  % Títulos y contenido
  \node[tit] at (0.1,\H-0.1) {Socios clave};
  \node[cue=4.9] at (0.1,\H-0.7) {\Socios};
  \node[tit] at (\c+0.1,\H-0.1) {Actividades clave};
  \node[cue=4.9] at (\c+0.1,\H-0.7) {\Actividades};
  \node[tit] at (\c+0.1,\m-0.1) {Recursos clave};
  \node[cue=4.9] at (\c+0.1,\m-0.7) {\Recursos};
  \node[tit] at (2*\c+0.1,\H-0.1) {Propuesta de valor};
  \node[cue=4.9] at (2*\c+0.1,\H-0.7) {\Propuesta};
  \node[tit] at (3*\c+0.1,\H-0.1) {Relación con clientes};
  \node[cue=4.9] at (3*\c+0.1,\H-0.7) {\Relacion};
  \node[tit] at (3*\c+0.1,\m-0.1) {Canales};
  \node[cue=4.9] at (3*\c+0.1,\m-0.7) {\Canales};
  \node[tit] at (4*\c+0.1,\H-0.1) {Segmentos de clientes};
  \node[cue=4.9] at (4*\c+0.1,\H-0.7) {\Segmentos};
  \node[tit] at (0.1,\h-0.1) {Estructura de costos};
  \node[cue=12.6] at (0.1,\h-0.7) {\Costos};
  \node[tit] at (\W/2+0.1,\h-0.1) {Fuentes de ingresos};
  \node[cue=12.6] at (\W/2+0.1,\h-0.7) {\Ingresos};
\end{tikzpicture}}

\vspace{2pt}
\noindent\scriptsize Basado en Osterwalder y Pigneur, \emph{Business Model Generation} (2010), CC BY-SA 3.0.
\val{H1} validada,\ \inval{H2} invalidada, (H3) pendiente. Evidencia en \texttt{hipotesis.csv} y \texttt{entrevistas.csv}.
\par\endgroup
. Este archivo no se compila solo.
% Llenar los nueve bloques. Cada afirmación termina con el identificador de la
% hipótesis que la respalda en hipotesis.csv, por ejemplo (H1).
% Marcar las hipótesis validadas con \val{H1} y las invalidadas con \inval{H2}.
\definecolor{navy}{HTML}{1F3467}
\providecommand{\val}[1]{\textcolor{green!50!black}{\textbf{(#1 \checkmark)}}}
\providecommand{\inval}[1]{\textcolor{red!70!black}{\textbf{(#1 $\times$)}}}
\providecommand{\pendiente}[1]{\textcolor{gray}{\emph{#1}}}

% ---------- Contenido de los bloques ----------
\newcommand{\Socios}{\begin{itemize}
  \item \pendiente{Proveedores y aliados sin los que el negocio no opera.} (H?)
\end{itemize}}
\newcommand{\Actividades}{\begin{itemize}
  \item \pendiente{Lo que el equipo hace todos los días para entregar valor.} (H?)
\end{itemize}}
\newcommand{\Recursos}{\begin{itemize}
  \item \pendiente{Datos, modelos, personas e infraestructura.} (H?)
\end{itemize}}
\newcommand{\Propuesta}{\begin{itemize}
  \item \pendiente{Qué dolor elimina o qué ganancia crea, para quién.} (H?)
\end{itemize}}
\newcommand{\Relacion}{\begin{itemize}
  \item \pendiente{Autoservicio, asistencia personal, comunidad.} (H?)
\end{itemize}}
\newcommand{\Canales}{\begin{itemize}
  \item \pendiente{Cómo se entera, compra y recibe el producto el cliente.} (H?)
\end{itemize}}
\newcommand{\Segmentos}{\begin{itemize}
  \item \pendiente{Perfil del cliente que paga y del usuario, si son distintos.} (H?)
\end{itemize}}
\newcommand{\Costos}{\begin{itemize}
  \item \pendiente{Costos fijos y variables principales (ver modelo financiero).} (H?)
\end{itemize}}
\newcommand{\Ingresos}{\begin{itemize}
  \item \pendiente{Qué se cobra, cuánto y con qué frecuencia.} (H?)
\end{itemize}}

% ---------- Dibujo. No editar debajo de esta línea ----------
\begingroup
\fontsize{9}{11}\selectfont
\setlist{nosep,leftmargin=1em}
\noindent{\Large\bfseries\color{navy} Business Model Canvas}\hfill
{\large \pendiente{Nombre de la empresa}}\hfill \pendiente{Versión y fecha}

\vspace{4pt}
\noindent
\resizebox{\linewidth}{!}{%
\begin{tikzpicture}[x=1cm,y=1cm]
  \def\W{25.8}\def\H{16.8}\def\h{5.0}
  \tikzset{caja/.style={draw=navy,line width=0.8pt},
           tit/.style={anchor=north west,font=\bfseries\color{navy}},
           cue/.style={anchor=north west,text width=#1cm,align=left}}
  % Retícula: 5 columnas arriba, 2 abajo
  \def\c{5.16}
  \draw[caja] (0,0) rectangle (\W,\H);
  \draw[caja] (0,\h) -- (\W,\h);
  \foreach \i in {1,2,3,4} \draw[caja] (\i*\c,\h) -- (\i*\c,\H);
  \draw[caja] (\W/2,0) -- (\W/2,\h);
  \pgfmathsetmacro\m{(\h+\H)/2}
  \draw[caja] (\c,\m) -- (2*\c,\m);
  \draw[caja] (3*\c,\m) -- (4*\c,\m);
  % Títulos y contenido
  \node[tit] at (0.1,\H-0.1) {Socios clave};
  \node[cue=4.9] at (0.1,\H-0.7) {\Socios};
  \node[tit] at (\c+0.1,\H-0.1) {Actividades clave};
  \node[cue=4.9] at (\c+0.1,\H-0.7) {\Actividades};
  \node[tit] at (\c+0.1,\m-0.1) {Recursos clave};
  \node[cue=4.9] at (\c+0.1,\m-0.7) {\Recursos};
  \node[tit] at (2*\c+0.1,\H-0.1) {Propuesta de valor};
  \node[cue=4.9] at (2*\c+0.1,\H-0.7) {\Propuesta};
  \node[tit] at (3*\c+0.1,\H-0.1) {Relación con clientes};
  \node[cue=4.9] at (3*\c+0.1,\H-0.7) {\Relacion};
  \node[tit] at (3*\c+0.1,\m-0.1) {Canales};
  \node[cue=4.9] at (3*\c+0.1,\m-0.7) {\Canales};
  \node[tit] at (4*\c+0.1,\H-0.1) {Segmentos de clientes};
  \node[cue=4.9] at (4*\c+0.1,\H-0.7) {\Segmentos};
  \node[tit] at (0.1,\h-0.1) {Estructura de costos};
  \node[cue=12.6] at (0.1,\h-0.7) {\Costos};
  \node[tit] at (\W/2+0.1,\h-0.1) {Fuentes de ingresos};
  \node[cue=12.6] at (\W/2+0.1,\h-0.7) {\Ingresos};
\end{tikzpicture}}

\vspace{2pt}
\noindent\scriptsize Basado en Osterwalder y Pigneur, \emph{Business Model Generation} (2010), CC BY-SA 3.0.
\val{H1} validada,\ \inval{H2} invalidada, (H3) pendiente. Evidencia en \texttt{hipotesis.csv} y \texttt{entrevistas.csv}.
\par\endgroup
